\chapter*{Abstract}
\addcontentsline{toc}{chapter}{Abstract}

\begin{english}
Energy renovation of the building stock is a central component of European
and national energy policy, as buildings account for a substantial share of
final energy consumption and greenhouse gas emissions. Although public
subsidy schemes, banking products and private financing mechanisms are
available, information about eligibility rules, programme terms and key
financial characteristics is often dispersed across websites, guides,
announcements and documents. This fragmentation makes it difficult for end
users, engineers and administrators to identify, compare and update relevant
information in a systematic way.

The objective of this diploma thesis is to develop a methodology and an
application for the automatic extraction and standardization of data related
to financing instruments for energy renovations in Greece. The proposed
approach combines web data collection techniques with large language models in
order to support source relevance assessment, structured information
extraction, storage in a common schema and quality control of the extracted
results. Emphasis is placed on preserving the connection between extracted
fields and primary sources, avoiding unsupported conclusions and maintaining
human oversight in a domain with financial and administrative implications.

The thesis aims to contribute to the reduction of information fragmentation
and to the more effective use of available financing instruments. At the same
time, it examines the capabilities and limitations of large language models in
an applied energy policy context where accuracy, traceability and continuous
updating are essential requirements.

\vspace{0.8cm}
\noindent\textbf{Keywords:} building energy renovation, financing instruments,
large language models, web scraping, information extraction, structured data,
energy policy.
\end{english}
\clearpage
